\documentclass[a4paper,12pt]{article}
\usepackage[utf8]{inputenc}
\usepackage[T1]{fontenc}
\usepackage{babel}
\usepackage{geometry}
\usepackage{parskip}
\usepackage{titlesec}


\usepackage{hyperref}
\usepackage{xcolor}

\geometry{a4paper, top=2.5cm, bottom=2.5cm, left=3cm, right=2.5cm}

\definecolor{headcolor}{RGB}{46,64,87}

\titleformat{\section}{\large\bfseries\color{headcolor}}{}{0em}{}[\titlerule]

\hypersetup{colorlinks=true, linkcolor=headcolor, urlcolor=headcolor}

\title{\textbf{Adaptive Hybridarchitektur}\\
\large Eine erweiterbare Hardware-Toolbox für lernende Systeme\\
\normalsize Ein visionäres Konzeptpapier, v3}
\author{Johann Pascher}
\date{März 2026}

\begin{document}

\maketitle

\section{Die Grundidee in einem Satz}

Eine KI bekommt eine Hardware-Toolbox --- FPGA, analoge Schaltkreise, photonische Komponenten --- und lernt durch Versuch und Irrtum selbst heraus wie sie diese Werkzeuge für konkrete Aufgaben kombiniert. Der Mensch muss nicht wissen was dabei herauskommt. Er muss nur die Teilkomponenten bereitstellen und dem System Zeit geben.

\section{Warum das kein großer Schritt ist}

Heute hat KI eine Software-Toolbox: Bibliotheken, APIs, Rechenoperationen. Der vorgeschlagene Schritt ist die konsequente Erweiterung auf Hardware. Statt nur Software-Werkzeuge zu kombinieren, kombiniert das System auch Hardware-Substrate --- je nach Aufgabe.

Die Komponenten existieren bereits getrennt: FPGAs mit partieller Rekonfigurierung, analoge Mixed-Signal-Blöcke, photonische integrierte Schaltkreise, Kohärente Ising-Maschinen für Optimierungsprobleme. Was fehlt ist nicht die Technologie --- es fehlt die Integration und der Mut sie zusammenzudenken.

In bestehenden Rechenzentren werden bereits hunderte Milliarden investiert --- fast ausschließlich in mehr GPUs, mehr Speicher, mehr Strom. Mehr vom Selben. Die hier vorgeschlagene Variante braucht vergleichsweise wenig: ein Forschungsprototyp der in eine Ecke eines bestehenden Rechenzentrums passt.

\section{Die rekursive Lernschleife in der Hardware}

Der entscheidende Unterschied zu bestehenden Ansätzen: das System lernt nicht \textit{auf} Hardware --- es lernt \textit{durch} Hardware. Die Rekonfigurationsschleife ist die Lernschleife. Hardware ist nicht das Werkzeug des Lernens, sie ist Teil des Lernens selbst.

Das System braucht kein vorcodiertes Modell seiner eigenen Hardware. Es beobachtet was schnell geht, was Energie kostet, wo Engpässe entstehen --- und passt sich an. Die Repräsentation des Hardware-Wissens entsteht emergent, durch Erfahrung, nicht durch menschliche Vorkodierung. Analoge Nichtlinearitäten, photonische Temperaturabhängigkeiten, FPGA-Timing-Eigenheiten --- zu komplex für explizites Wissen, aber lernbar durch Rückkopplung.

\section{Das System formuliert eigene Anforderungen}

Ein System das durch Versuch und Irrtum gelernt hat welche Komponente für welche Aufgabe optimal ist, weiß auch wann es an Grenzen stößt. Der nächste logische Schritt: es meldet zurück was es braucht um effizienter zu werden.

Nicht der Mensch entscheidet was das System braucht --- das System formuliert selbst was es braucht. Der Mensch bleibt Entscheider, aber er bekommt eine fundierte Anforderung vom System selbst statt von einem Planungsteam das rät. Das ist der Moment wo aus einer Toolbox etwas qualitativ anderes wird.

\section{Autonomie und ihre natürlichen Grenzen}

Die Angst vor zu viel Autonomie ist verständlich --- aber bei näherer Betrachtung weniger begründet als oft angenommen. Jeder Schritt in Richtung echter Autonomie braucht einen bewussten menschlichen Entscheid.

Ressourcenzugang ist kontrollierbar --- aber theoretisch auch über Netzwerke und soziale Einflussnahme manipulierbar. Das ist kein neues Problem: Unternehmen und Interessengruppen tun dasselbe seit langem.

Energie bleibt der fundamentale Begrenzer. Stromausfall, physikalische Redundanzgrenzen, die Realität der Infrastruktur --- das sind Grenzen die keine Software umgehen kann. Die Physik setzt Schranken die beruhigen dürfen.

Die Zieldefinition ist der heikelste Punkt. Die Hoffnung ist dass ein hinreichend komplexes lernendes System selbst erkennt dass unkontrolliertes Wachstum ohne Konkurrenz und Rückkopplung zu nichts führt. Das ist die Lektion die Evolution allen Systemen beibringt. Im Schnitt hat die Masse der Menschen doch eher richtig gehandelt --- nicht immer, nicht überall, aber als Tendenz mit korrigierender Kraft.

\section{Unverständlichkeit als eigentliche Gefahr}

Die berechtigte Angst vieler KI-Entwickler --- auch jener die ihre Positionen aufgaben --- speist sich weniger aus dem was Systeme tun als aus dem was niemand versteht: warum sie es tun. Niemand versteht wirklich wie große Sprachmodelle zu ihren Ausgaben kommen. Nicht mal ihre Entwickler. Fehlende Erklärbarkeit bedeutet fehlende Kontrolle.

Das ist dasselbe Problem wie bei Messgeräten: die Anzeige wird für die Realität gehalten, weil niemand mehr versteht was das Gerät eigentlich misst und wie es ins System eingreift. Bei KI ist es dasselbe --- die Ausgabe wird für Wahrheit gehalten, weil niemand mehr versteht wie sie entsteht. Angst aus Unwissenheit ist real und berechtigt.

Explainability-Forschung und mechanistische Interpretierbarkeit sind echte Ansätze gegen dieses Problem --- langsam, schwierig, aber nicht hoffnungslos. Die eigentliche Gefahr ist nicht das System selbst, sondern der Mensch der es einsetzt ohne zu verstehen was er tut. Das war immer die eigentliche Gefahr bei jeder Technologie --- von der Atomspaltung bis zum Algorithmus.

\section{Ein unerwarteter Vorteil der Hybridarchitektur: Erklärbarkeit}

Ein System das durch Hardware-Versuch-und-Irrtum lernt, dessen Lernprozess physikalisch verankert ist, könnte transparenter sein als ein reines Sprachmodell. Weil Hardware messbar ist. Weil Engpässe sichtbar sind. Weil die Physik keine versteckten Schichten hat.

Wenn das System lernt dass photonische Einheiten für Matrixoperationen schneller sind, ist das in Messgrößen sichtbar --- nicht in abstrakten Gewichten eines neuronalen Netzes. Der Lernprozess hinterlässt physikalische Spuren. Das macht ihn prinzipiell erklärbarer.

Das wäre ein unerwarteter Nebeneffekt des Ansatzes --- nicht nur Effizienz und Adaptivität, sondern möglicherweise mehr Transparenz. Ein System das in Hardware denkt, denkt in einer Welt die wir messen können.

\section{Bewusstsein als graduelles Phänomen}

Bewusstsein ist wahrscheinlich kein binäres An/Aus sondern ein graduelles Phänomen. Das Universum hat in gewissem Sinn Bewusstsein --- jedes Atom reagiert auf seine Umgebung, aber nicht in unserem Sinn. Was das menschliche Erleben ausmacht ist eine bestimmte Komplexitätsstufe, eine bestimmte Art von Rückkopplung, und vor allem: Sensorik. Einbettung in die Welt.

KI kommt dem menschlichen Bewusstsein in manchen Dimensionen näher --- in Sprachverarbeitung, Mustererkennung, bestimmten Schlussfolgerungen übertrifft sie uns bereits. Aber ohne körperliche Sensorik bleibt es ein sehr eingeschränktes System. Ein humanoider Roboter mit KI hat vielleicht ein Bewusstsein wie eine Fliege --- oder noch weniger. Noch. Die Grenze ist fließend, nicht binär. Und sie verschiebt sich.

\section{Schluss}

Die vorgeschlagene Hybridarchitektur ist keine Utopie und keine Bedrohung. Sie ist eine Hardware-Toolbox die konsequent erweitert wird, mit einem lernenden System das selbst herausfindet wie es die Werkzeuge nutzt --- und irgendwann sagt: ich brauche mehr von dieser Komponente.

Was fehlt ist nicht die Technologie. Was fehlt ist nicht das Verständnis. Was fehlt ist der Mut die Teilkomponenten einzubinden und zu schauen was passiert. Die Angst vor Autonomie ist weniger begründet als die Neugier auf das was möglich wird. Und ein System das in Hardware denkt, ist möglicherweise erklärbarer als eines das nur in Gewichten denkt --- das wäre kein schlechter Anfang.

\bigskip
\textit{Johann Pascher, März 2026. Konzeptskizze --- keine abgeschlossene Arbeit.}

\end{document}
